\subsection{Historical CCME 1.0.0 obligations}
\label{sec:obligations-historical}

Formalising CCME 1.0.0 at \texttt{713f1a0b} exposed six failure classes. Each obligation below
states the claim, a minimal counterexample, the consequences, a repair, and a
conformance vector. Three concern \spec{13.7c}, one \spec{13.7b}, one
\spec{9.3}, and one the optimisation advice of \spec{13.11}.

None is a defect in the reference implementation. In two cases
(\cref{ob:pre,ob:duality}), the implementation already uses a mechanism absent
from the specification. A second implementation following only the written
contract would therefore be incorrect. Source references below are
relative to \code{services/dispat/internal/plan/} and cite the line numbers of
the released \code{services/dispat/v1.7.2} tag. Those sources have since
moved: at \code{252216ba}, \code{containedInBaseline} is at \code{plan.go}
line 2507, \code{Plan.StaleSources} at line 1280, and the per-unit traversal
of \code{propagate.go} is served from a walk cache (line 76) whose
composition is checked against the literal walk (\cref{sec:conformance}).

Released CCME 2.0.0 resolves them as follows: \cref{ob:pre} by $W_{\rm fresh}$;
\cref{ob:g5,ob:origin} by the full retry invariant and fixed-state wording;
\cref{ob:duality} by applying channel eligibility in the upward audit;
\cref{ob:g8} by stating axis independence in the bump-to-channel direction
only; and \cref{ob:bucket} by retaining each unit's own source exclusion when
bucketed. The counterexamples remain regression evidence for those normative
conditions. The regression vectors named below as ``the released 2.0.0
form'', 133 to 138, are in the current text the edge cases 133 to 138 of
\spec{15.7}, one row each. Except where a paragraph explicitly names CCME 2.0.0 or a current
theorem, every numbered \spec{} reference and quoted guarantee in this subsection
refers to the CCME 1.0.0 snapshot at \code{713f1a0b}; present tense describes
the historical rule being executed, not the released 2.0.0 text.

\subsubsection{The prerelease window never discharges}

\begin{obligation}[the 1.0.0 G4/G6 claims fail on a prerelease train]
\label{ob:pre}
Under the 1.0.0 single-window rule, ledger discharge and convergence fail after a successful prerelease publish. Current \cref{thm:g4,thm:g6} use $W_{\rm fresh}$ and do not make that claim.
\end{obligation}

\paragraph{Counterexample (bump axis).}
This is \SPEC{} vector 95b, run twice. Baselines
\code{core@1.5.0-beta.0} and \code{cli@2.0.1-beta.0}, both with an older stable
baseline; \code{cli} depends on \code{core}; one commit
\code{fix(core)\^{}: x}.

Run 1 plans \code{core@1.5.0-beta.1} and \code{cli@2.0.1-beta.1} and publishes
both. Run 2, at the same \Head{}: by \cref{lem:window} neither window changed,
because \cref{def:window} measures from the last \emph{stable} tag and no
stable tag was written. Under the 1.0.0 direct and propagated rules, \code{core} and \code{cli} each receive \bpatch{} again, both packages are planned again, and
\spec{11.4} increments the counters to \code{beta.2}. Run 3 gives
\code{beta.3}. The plan never empties, and \diag{E199}, the convergence
check of \spec{19.6}, which admits no exclusion but held packages, fires
after the first fully successful run of any prerelease train.

\paragraph{Counterexample (channel axis).}
\Cref{thm:g7} says the channel a package arrived on is not proposed for it
again. It does not say nothing is proposed. Let $p$ have stable baseline
\code{1.2.0} and a window containing an older \code{release(p)\%rc} and a
newer \code{release(p)\%beta}. Run 1: \spec{11.6} precedence gives the newer
directive, $p$ is released \code{1.2.1-beta.0} (\diag{W186} noting the
conflict), and the operator's stated intent is honoured. Run 2: the window is
unchanged, so both directives are still candidates; \code{\%beta} now proposes
nothing (\diag{W199}, \cref{thm:g7}) and the \emph{older} \code{\%rc},
which \spec{11.6} had ruled had lost, wins by default and $p$ is released
\code{1.2.1-rc.0}. Run 3: \code{\%rc} is now inert and \code{\%beta} proposes
\code{beta} again, computing \code{1.2.1-beta.0}, which does not exceed the
baseline \code{1.2.1-rc.0}, so the run stops with \diag{E195}. A repository
that did nothing wrong is wedged, having published a version onto a channel its
newest directive did not ask for.

\paragraph{What broke in 1.0.0.} Its G4 (no double release), G6 (convergence), and G7 channel-settling claim (the channel axis settles) all fail, and with
them \spec{17.2}'s idempotency clause and
\spec{19.6}'s convergence check, for every package on a prerelease train,
which is to say for the entire feature the prerelease flow exists to support.

\paragraph{The repair.} The window answers one question, what has $p$ not
yet released, and \spec{13.3} correctly measures it from the stable tag
because \spec{11.4} recomputes a train's target over the whole train. The
missing distinction is that the same set is then used to answer a second
question: is anything still \emph{owed}? Separate them. Define the
\emph{fresh window}
\[
  \Winf{\st}{p} \;=\; \Win{\st}{p} \setminus \Rch{\basec{\st}{p}},
\]
which for a package whose baseline is stable equals $\Win{\st}{p}$, and for one
on a train excludes the commits its own prereleases already published. Then:
\begin{enumerate}[label=(R\arabic*),leftmargin=3.2em]
  \item \cref{def:sources} discharges on $\Winf{\st}{p}$, so a package whose
    prerelease shipped the work counts as a source (which \spec{13.4a} already
    intends: the artefact is public);
  \item plan membership (\cref{def:plan}) requires a contribution from a commit
    in $\Winf{\st}{p}$, while the version computation (\spec{11.4}) keeps
    reading $\Win{\st}{p}$, so the train's target still accumulates over the
    whole train;
  \item candidate collection in \cref{def:phase1,def:channels} is restricted to
    $\Winf{\st}{p}$, which is what stops the superseded directive of the second
    counterexample from voting again.
\end{enumerate}
With (R1)--(R3), \cref{lem:window} strengthens to $\Winf{\st'}{p} = \emptyset$
for \emph{every} released $p$, prerelease or not, because $\basec{\st'}{p} =
\Head$. Fresh discharge repairs the prerelease counterexample. Current \cref{thm:g4,thm:g5,thm:g6} separately state the full retry-invariant conditions needed for persistence, non-widening, and bounded convergence; (R1)--(R3) alone do not imply them.

\paragraph{Evidence.} The reference implementation has exactly this predicate,
under a different name and undocumented in \SPEC{}:
\code{containedInBaseline} (\code{plan.go},
lines 1836--1855) reports whether a commit has already been published by the
package's baseline tag, and its callers implement (R1) in
\code{sourcePackages}, (R2) by gating \code{NewWork} on it, and (R3) in both
\code{resolveChannels} and \code{propagateChannels}. Its comment states the
consequence the specification does not: without it the engine ``would re-warn
\dots on every run until graduation''.

\paragraph{Regression vectors.} Vector 133 is the released 2.0.0 form of the bump-axis case. Historically, (i) vector 95b, run twice: the second run's plan
must be empty. (ii) The two-directive window above, run three times: run 2 must
be empty, and \diag{E195} must not arise.

\subsubsection{Window-scoped suppression lapses, so the admitted set can grow}

\begin{obligation}[the 1.0.0 G5 claim lacks suppression stability]
\label{ob:g5}
The 1.0.0 proof assumed ``same sources, same depth, same graph at \Head{}'', although window-scoped suppression can change the source set. CCME 2.0.0 makes the needed stability part of the full retry invariant used by current \cref{thm:g5}.
\end{obligation}

The root cause is a design decision taken three times, in three sections, each
time for a good local reason: suppression is scoped to the window. A
\code{cancel} stops counting once the package releases (\spec{13.4a}); a hold
does the same (\spec{13.4a}); and a correction ``reaches only undischarged
work'', so a \code{Deletes} is a no-op for a package that has released its
target (\spec{7.4.2}, \cref{rem:corrections}). Each rule is right on its own
terms. Together they mean that the set of units a package acts as a propagation
source for is \emph{monotone in the state}, which is precisely what \G{5} says
it is not.

\paragraph{Counterexample (cancellation).} \code{cli} depends on \code{core}.
Three commits: $A$ carrying \code{feat(core)\^{}: streaming reader}, then $B$
carrying \code{cancel(core)}, then $C$ carrying \code{fix(core): unrelated}, in
that ancestry.

Run 1: $\canc{A}{\code{core}}$ holds and $A \in \Win{\st_1}{\code{core}}$, so
\code{core} is not in $\srcp{\st_1}{u_A}$ and $u_A$ propagates to nothing;
\code{core} releases \code{1.4.3} from $C$ alone (\SPEC{} vector 41).
$\plan{\st_1} = \{\code{core}\}$: \code{cli} is not planned.

Run 2: \code{core} is tagged at \Head{}, so $\Win{\st_2}{\code{core}} =
\emptyset$, so $\disch{\st_2}{\code{core}}{A}$ holds, so \code{core} re-enters
$\srcp{\st_2}{u_A}$ by the discharge disjunct of \cref{def:sources}. Now $u_A$
propagates: $A \in \Win{\st_2}{\code{cli}}$, \code{cli} is not cancelled, and
the bump is admitted. $\plan{\st_2} = \{\code{cli}\} \not\subseteq
\plan{\st_1}$.

\paragraph{Counterexample (corrections).} The same shape without any
\code{cancel}. Let $B$ instead carry \code{chore(core): drop it} with
\code{Deletes: <A>}. In run 1 the correction discards $A$'s record for
\code{core}, so $u_A$ has no source and reaches nothing; \code{core} still
releases, on $C$. In run 2, $\Win{\st_2}{\code{core}}$ is empty, so by
\spec{7.4.2} the correction is a discharged no-op for \code{core}
(\diag{W209}), $A$'s record stands again, and $u_A$ propagates to \code{cli}
exactly as before. This route is worth stating separately because \diag{W209}
reads as reassurance (the correction did not need to do anything) when it
is in fact the moment the suppressed propagation returns.

\paragraph{What broke in 1.0.0.} Its unrestricted G5 claim and the ``nothing is added'' step of its G6 proof failed. \Cref{thm:g6}'s conclusion survives here (\code{cli}
releases in run 2 and run 3 is empty) but the variant argument does not, and
the security reading of \G{5} does: whether \code{cli} is released at all
depends on whether \code{core}'s publish succeeded. Inducing a publish failure
changes the set of packages a commit eventually releases, which is the
inference \spec{18.1} relies on \G{5} to forbid.

\paragraph{The repair, and why the obvious one is wrong.} The tempting fix is
to make cancellation sticky, never letting a cancelled package become a source,
by defining $\disch{\st}{p}{c} \land \lnot\canc{c}{p}$. It is wrong: it
breaks \SPEC{} vector 72, where \code{cancel(core)} lands \emph{after}
\code{core} published the work and the consumers' catch-up must still proceed.
Ancestry cannot separate the two cases, since in both the cancel is in the
history at \Head{} and in both \code{core} carries a tag at \Head{}. The
difference is only whether the released version counted the unit, and that is
not recoverable from the state.

The recommendation is therefore to restate the guarantee rather than change the
algorithm: \G{5} holds for every unit none of whose sources is suppressed at
\Head{}, which is hypothesis (H2), and the widening in the remaining case is
intended behaviour, since neither a \code{cancel} nor a \code{Deletes} touches
code (\spec{10.1}, \spec{7.4.2}) and the work of $A$ is in fact inside the
published \code{core@1.4.3}. What the specification should not do is claim a
subset property that the discharge rule of \spec{13.4a} deliberately violates,
and \spec{18.1} should not rest a security argument on the unrestricted form.

\paragraph{Evidence.} The cancellation counterexample was run against the
reference implementation, which follows \spec{13.4a} literally and reproduces
it: with the history above and \code{core} tagged at \Head{} in the second
state, the first plan releases \code{core} alone and the second releases the
consumer as well, reporting it as a \emph{catch-up from core}. The engine is
doing what the specification says; it is \G{5} that does not describe it.

\paragraph{Regression vector.} Vector 134 is the released 2.0.0 form. In the historical pair, each history above is run twice: run 2
plans \code{cli}, the plan is documented as a superset of run 1's, and the
\diag{W193} naming \code{core} as the origin is present so that the widening is
visible to whoever reads the plan rather than silent.

\subsubsection{The staleness dual omits the resolvability test}

\begin{obligation}[\spec{13.7b}'s two formulations disagree]
\label{ob:duality}
\spec{13.7b} states that \code{staleSources(D)} ``is non-empty exactly when
\S9.2 assigns \code{D} a non-\code{none} propagated bump'', and that
disagreement is an implementation bug. The procedure tests scope, depth,
window, and cancellation, but has no counterpart to \cref{a:resolvable}.
\end{obligation}

\paragraph{Counterexample.} \SPEC{} vector 95 (also edge case 60m):
\code{feat(core)\^{}\%beta: x} where \code{core} moves onto \code{beta} and its
consumers stay on \code{stable}. \spec{9.3a} suppresses every consumer's bump
with \diag{W208}, so $\propb{\st}{\code{ui}} = \bnone$, and \SPEC{} is
explicit that ``\code{feat(core)\^{}\%beta} releases \code{core} alone''. But
\code{staleSources(ui)} passes every guard it has and returns the row
$(\code{core}, u, 1, \bpatch)$. The audit reports \code{ui} as behind a
dependency for which no release is owed, and by \spec{13.7b}'s own words the
implementation is now buggy by definition.

\paragraph{The repair.} Add the hoisted predicate of \cref{prop:hoist} to the
procedure, immediately after the \code{cancelledFor} test:
\begin{speccode}
        srcChan = { channel[P] : P in sources }        # §9.3a, as in §9.2
        if not ('stable' in srcChan or channel[D] in srcChan):  continue
\end{speccode}
With it, \cref{thm:duality} holds unconditionally and its hypothesis can be
dropped.

\paragraph{Evidence.} The reference implementation cannot exhibit the
disagreement, because \code{Plan.StaleSources}
(\code{plan.go}, line 1217) returns the rows
recorded by the downward traversal rather than re-deriving them from an
independent upward walk. The duality is therefore true there by construction
and untested; an implementation that follows \spec{13.7b} literally, as a
second implementation would, disagrees on every prerelease workspace.

\paragraph{Regression vector.} Vector 135 is the released 2.0.0 form; historically, vector 95 was extended: assert that
\code{staleSources(ui)} is empty, and that the randomised
down-versus-up cross-check \spec{13.7b} recommends is run on a workspace
containing at least one prerelease train.

\subsubsection{Axis independence is one-directional}

\begin{obligation}[the ``vice versa'' in \G{8}]
\label{ob:g8}
\spec{13.7c} states \G{8} as ``suppressing a bump under \spec{9.3a} never
suppresses a channel, and vice versa'', and proves it by observing that the two
passes share only $\chanres{\st}{\cdot}$, which phase 3 reads and never writes. That
argument proves one direction only, and the other is false.
\end{obligation}

\paragraph{Counterexample.} Let \code{core} and its consumer \code{app} both
sit on \stable, and let one commit carry \code{feat(core)\^{}\%beta++1}: the
caret propagates a bump to \code{app} and the \code{++1} propagates the
\code{beta} channel to it, which is exactly the pairing \spec{9.3a} prescribes
for taking consumers along. Suppose \code{app}'s own window also contains
\code{release(app)\%rc}. By \cref{def:channels} a direct directive beats every
propagated one, so $\chanres{\st}{\code{app}} = \code{rc}$ and the propagated
\code{beta} is discarded. Phase 3 then computes
$\srcch{\st}{u} = \{\code{beta}\}$, finds
$\chanres{\st}{\code{app}} = \code{rc} \notin \srcch{\st}{u}$ and no source on
\stable, and suppresses the bump with \diag{W208}. Had the propagated channel
been applied, \code{app} would have been on \code{beta} and the bump admitted.
The suppression of a channel proposal caused the suppression of a bump.

\paragraph{The repair.} State \G{8} as \cref{thm:g8} does: the channel
assignment is computed by phases 1--2 and reads no phase-3 output, so no bump
suppression can affect it. The dependency in the other direction is real, is
the entire content of \spec{9.3a}, and should be documented as such rather than
denied. The operational consequences \spec{13.7c} draws from \G{8} are
unaffected: they only ever use the direction that holds.

\paragraph{Regression vector.} Vector 136 is the released 2.0.0 form. The historical workspace above asserts that \code{app} is
released on \code{rc} with no propagated bump and a \diag{W208}, and that the
\diag{W208} is a consequence of the channel precedence rule rather than of
anything the caret did.

\subsubsection{Inherited channels read a baseline that release moves}

\begin{obligation}[\spec{9.3}'s replay-independence claim]
\label{ob:origin}
\spec{9.3} states of the \code{inherit} resolution that ``the result is
deterministic, does not depend on which earlier run published the origin''.
The second half is false: resolution rule 2 reads
$\chanof{\base{\st}{q}}$, and the origin's own release moves it.
\end{obligation}

\paragraph{Counterexample.} \code{app} depends on \code{core}; \code{core}
sits at baseline \code{1.5.0-beta.0} over stable baseline \code{1.4.2};
\code{app} at \code{2.0.0}. Two commits: $C_1$ carrying
\code{feat(core)\^{}++1: x} (no explicit channel value, so the propagated
channel is \code{inherit}) and $C_2$ carrying
\code{release(core)\%stable: graduate}.

Run 1: \code{core}'s own channel is \stable{} (the direct directive), but
the inherited value for \code{app} is
$\chanof{\base{\st_1}{\code{core}}} = \code{beta}$, so \code{app} is
planned at \code{2.0.1-beta.0} while its provider graduates to
\code{1.5.0}. Suppose \code{core} publishes and \code{app} fails. Run 2:
$\chanof{\base{\st_2}{\code{core}}} = \stable$, the proposal is inert
(\diag{W199}), and \code{app} is planned at \code{2.0.1}. Both runs plan
\code{app}, at different versions, with every completed release stable:
the 1.0.0 cross-run version-stability claim fails. CCME 2.0.0 instead fixes inherited baseline-channel state as one component of its full retry invariant, as required by current \cref{thm:g3}. Run against the reference implementation, the two
plans are exactly as computed here (\code{app}: \code{beta}
$\rightarrow$ \code{stable}, \code{2.0.1-beta.0} $\rightarrow$
\code{2.0.1}).

\paragraph{What breaks, and what does not.} Determinism
(\cref{thm:determinism}) is unaffected: each plan is still a pure function
of its state. What fails is the \emph{cross-run} half of the specification's
sentence, and with it channel stability for unreleased dependents of a
graduating origin. Note also that run 1's own output is operationally odd:
a consumer is dragged onto a train in the same run that ends it. Both
effects have one cause, phase 1 reading baselines so that it cannot depend
on phase 2's output (\spec{9.2}); the price of that layering is that
\code{inherit} answers with the origin's pre-graduation channel.

\paragraph{The repair.} Scope the specification's sentence to a fixed state
(determinism per run, which is what \cref{thm:determinism} proves), and
state the drift explicitly. Resolving \code{inherit} against the origin's
\emph{resolved} channel instead would remove the drift but requires a
fourth phase, which \spec{9.2}'s no-circularity argument was written to
avoid; the fresh window of \cref{ob:pre} does not remove it either, since
the commit remains pending for the failed dependent. The honest fix is
documentation plus the vector below.

\paragraph{Regression vector.} Vector 137 is the released 2.0.0 form: the two commits above, two runs: run 2 must
plan \code{app} at \code{2.0.1} on \stable{}, and the plan should carry the
\diag{W199} that records why the run-1 proposal lapsed.

\subsubsection{Bucketed traversal over-admits a unit's own sources}

\begin{obligation}[the 1.0.0 bucketing remedy has its own edge]
\label{ob:bucket}
CCME 1.0.0 \spec{13.11} warned that subtracting the merged source set from a bucketed
traversal silently withholds bumps, and prescribes admitting ``any node reached
across at least one edge within the depth bound''. By \cref{prop:bucket} that
prescription over-admits the set $X$: a package that is itself a source of the
unit that reaches it.
\end{obligation}

\paragraph{Counterexample.} One unit, \code{fix(core,cli)\^{}major: x}, where
\code{cli} depends on \code{core}. \spec{9.2} seeds \code{seen} with both
sources, so \code{cli} is not a propagation target and takes the unit's own
\bpatch. Under the merged rule, \code{cli} is reached from \code{core} across
one edge inside the depth bound and takes the propagated \bmajor{} as well:
a two-level difference in the published version, with no diagnostic.

\paragraph{The repair.} Bucket the traversal, then re-decide the few targets
that are themselves sources. Since the bucket key fixes the propagated bump,
the only question for such a $d$ is whether \emph{some} unit in the bucket
admits it, that is whether $\exists i: d \notin S_i \land \dist{S_i}{d} \leq
\ell$. The candidates are confined to $S$, which is small, and the per-unit
source sets are already materialised, so the check is cheap. Equivalently, and
more simply, exclude from bucketing any unit whose own source set contains a
package reachable from that set within the depth bound.

\paragraph{Evidence.} The reference implementation does not bucket: it
traverses once per unit (\code{propagate.go},
line 401), so it is unaffected. The obligation is on the advice, which a
performance-minded second implementation would follow.

\paragraph{Regression vector.} Vector 138 is the released 2.0.0 form. Historically, \SPEC{} Appendix B.7a, whose vectors test
exactly this class of almost-equivalence, gains one: \code{fix(core,cli)\^{}major}
must publish \code{cli} as a patch, bucketed or not.
